\section{Sequence Models}

In the preceding chapter, we explored how recurrent architecture, specifically LSTM and GRU cells, were developed to alleviate 
the vanishing gradient problem and better capture long-term dependencies in sequential data. 
Despite these advancements, RNN-based models still encounter significant challenges as input sequences grow in length.

This limitation is particularly prominent in the \textbf{encoder-decoder} (or sequence-to-sequence) framework. 
In this setup, an encoder RNN compresses the entirety of the input sequence into a single, fixed-length context vector, which the decoder RNN then utilizes to generate the output. 
Forcing the full semantic meaning of a variable-length input into a solitary vector creates a severe \textbf{information bottleneck}: 
as the sequence length increases, it becomes exceedingly difficult for this vector to retain all relevant details.

To resolve this bottleneck, the \textbf{attention mechanism} was introduced well before the advent of the Transformer architecture. 
Rather than relying on a static summary vector, attention enables the decoder to access the complete sequence of encoder hidden states. 
This allows the model to dynamically weigh and focus on the most pertinent information at each step of the generation process.

In this chapter, we examine sequence models that retain the RNN as their foundational building block but augment it with attention. 
We explore this synthesis through two representative tasks: \textbf{image captioning} and \textbf{neural machine translation}. 
In the subsequent chapter, we will discuss how the attention mechanism ultimately rendered recurrence obsolete, paving the way for the Transformer architecture.

\subsection{Image Captioning}
Image captioning is the task of generating a textual description for a given image. This task combines \textbf{computer vision} for the image understanding 
and \textbf{natural language processing} for generating the caption.
The main challenge in image captioning is to effectively capture the relevant features of the image and generate a coherent and contextually appropriate caption.

The usual approach to image captioning involves using a convolutional neural network (CNN) to extract features from the image,
and then using a recurrent neural network (RNN) to generate the caption based on those features.

\subsubsection{No Attention - Show and Tell}
\textbf{Show and Tell} is one of the early models for image captioning. Following the encoder-decoder architecture proposed for machine translation, 
Show and Tell uses an architecture combining two main components:

\begin{itemize}
    
    \item \textbf{Encoder}: A CNN (such as Inception or ResNet) is used to extract a fixed-length feature vector from the input image. 
    This feature vector captures the high-level visual information of the image.
    
    \item \textbf{Decoder}: An RNN (such as LSTM) is used to generate the caption word by word. 
    The RNN takes the feature vector from the encoder as input and generates the caption sequentially, predicting one word at a time until it generates an end token.
\end{itemize}

\begin{figure}[H]
    \centering
    \resizebox{\textwidth}{!}{\begin{tikzpicture}[
    >=Stealth,
    rnn/.style={rectangle, rounded corners=3pt, draw, fill=rnnblock, thick,
                minimum width=1.3cm, minimum height=0.65cm},
    clf/.style={rectangle, rounded corners=3pt, draw, fill=rnnblock!75, thick,
                minimum width=1.3cm, minimum height=0.65cm, font=\small},
    word/.style={font=\small\itshape, text=dlgreen!70!black},
    arr/.style={->, thick}
]
    % ---- CNN encoder ----
    \node[trapezium, trapezium angle=65, draw, fill=rnnblock, thick,
          minimum width=1.0cm, minimum height=0.7cm, font=\small] (cnn) at (0,0) {CNN};

    % ---- columns ----
    \def\dx{2.0}
    % image cell (receives CNN), then 5 word cells
    \node[rnn] (h0) at (1.6,0) {};
    \foreach \i/\out/\in in {1/The/START, 2/kids/The, 3/are/kids, 4/jumping/are, 5/STOP/jumping} {
        \node[rnn] (h\i) at ({1.6+\i*\dx},0) {};
        \node[clf, above=0.85cm of h\i] (c\i) {Classifier};
        \node[word, above=0.5cm of c\i] (y\i) {``\out''};
        \node[word, below=0.9cm of h\i] (x\i) {``\in''};
        \draw[arr] (x\i.north) -- (h\i.south);
        \draw[arr] (h\i.north) -- (c\i.south);
        \draw[arr] (c\i.north) -- (y\i.south);
    }
    % CNN feeds the image cell
    \draw[arr] (cnn.north) -- (h0.south);

    % ---- recurrent connections with hidden-state labels ----
    \foreach \i [evaluate=\i as \j using int(\i+1)] in {0,1,2,3,4} {
        \draw[arr] (h\i.east) -- node[below, font=\small] {$h_{\i}$} (h\j.west);
    }
\end{tikzpicture}
}
    \caption{Show and Tell: a CNN encodes the image into a feature vector that initializes the RNN decoder, which generates the caption word by word.}
\end{figure}
At the first time step, the RNN receives as input the \textbf{feature vector} from the encoder and a special start token.
At each subsequent time step $t$, the RNN receives as input the word generated at time step $t-1$ and the hidden state from the previous time step.
The RNN continues to generate words until it produces an end token, indicating the end of the caption.

\paragraph{Limitations}
Despite its success, this model faced several inherent limitations:

\begin{itemize}
    \item \textbf{Information Bottleneck}: The fixed-length feature vector may not capture all relevant details, especially in complex images with multiple interacting objects.
    
    \item \textbf{Vanishing Visual Context}: Because the RNN only "sees" the image at the first time step, the visual influence can weaken as the sentence grows, leading the model to rely more on language probability than the actual image.
    
    \item \textbf{Static Features}: The feature vector is global and static; the model cannot "shift its gaze" to different parts of the image as it describes different objects.
\end{itemize}

\subsubsection{Inference}
At inference time, two different strategies can be used to generate captions:

\begin{itemize}
    
    \item \textbf{Sampling}: At each time step, the model samples a word from the predicted probability distribution.
    The corresponding word is then fed as input to the next time step. This process continues until an end token is generated or a maximum length is reached.
    While this can lead to more diverse and creative captions, it also increases the risk of generating grammatically incorrect or nonsensical sequences.
    
    \item \textbf{Beam Search}: Instead of sampling, beam search keeps track of the top $k$ most likely sequences at each time step, 
    where $k$ is the beam width. At each step, the model expands each of the $k$ sequences by one word and keeps only the top $k$ resulting sequences based on their cumulative probabilities.
    Beam search is more computationally expensive than sampling but often results in more coherent and higher-quality captions.
\end{itemize}

\subsubsection{Evaluation Metrics}
In order to evaluate the performance of image captioning models, we must be able to compare the generated captions with
human-generated reference captions. To address this, the BLEU (Bilingual Evaluation Understudy) metric was introduced.

BLEU is a precision-based metric that measures the overlap between the generated caption and one or more reference captions.
It calculates the \textbf{n-gram precision}, which is the proportion of n-grams (contiguous sequences of n words)
in the generated caption that are also present in the reference captions.


\subsubsection{Attention: Show, Attend and Tell}

To overcome the limitations of the baseline \textit{Show and Tell} architecture, the \textit{Show, Attend and Tell} model introduces an \textbf{attention} mechanism. 
Rather than compressing the entire image into a single fixed-length vector, attention enables the model to dynamically focus on specific spatial regions of the image as it generates each word of the caption.

For example, consider an image of a dog playing with a ball in a park. 
When predicting the word "dog," relying on the global image context is suboptimal due to irrelevant background noise. 
Instead, the attention mechanism allows the network to isolate and "attend" specifically to the pixels containing the dog. 
By aligning local visual features with the exact word being generated at each time step, this dynamic focusing yields more accurate and context-aware captions.

\begin{figure}[H]
    \centering
    \begin{tikzpicture}[
    >=Stealth,
    rnn/.style={rectangle, rounded corners=3pt, draw, fill=rnnblock, thick,
                minimum width=1.2cm, minimum height=0.65cm},
    ctx/.style={rectangle, draw, fill=attncontext!70, thick, minimum width=0.85cm, minimum height=0.6cm, font=\small},
    yin/.style={rectangle, draw, fill=dlgreen!18, thick, minimum width=0.85cm, minimum height=0.6cm, font=\small},
    pout/.style={rectangle, draw, fill=attnscore!45, thick, minimum width=0.95cm, minimum height=0.6cm, font=\small},
    lbl/.style={font=\small},
    arr/.style={->, thick},
    grid3/.pic={\fill[highlightfill!40] (-0.3,-0.3) rectangle (0.3,0.3);
                \draw[step=0.2,thin] (-0.3,-0.3) grid (0.3,0.3);
                \draw[thin] (-0.3,-0.3) rectangle (0.3,0.3);}
]
    % ---- CNN + feature map ----
    \node[trapezium, trapezium angle=65, draw, fill=rnnblock, thick,
          minimum width=0.9cm, minimum height=0.6cm, font=\small] (cnn) at (0,-1.25) {CNN};
    \pic at (0,0) {grid3};
    \node[lbl, align=center] at (-1.15,0) {feature\\ map};
    \draw[arr] (cnn.north) -- (0,-0.32);

    % ---- decoder RNN cells ----
    \node[rnn] (h0) at (2,0) {$h_0$};
    \node[rnn] (h1) at (4,0) {$h_1$};
    \node[rnn] (h2) at (6,0) {$h_2$};
    \draw[arr] (0.32,0) -- (h0.west);
    \draw[arr] (h0) -- (h1);
    \draw[arr] (h1) -- (h2);
    \draw[arr] (h2.east) -- ++(0.9,0);

    % ---- attention maps above each cell ----
    \foreach \i/\a in {0/1,1/2,2/3} {
        \pic at ({2+\i*2},2.0) {grid3};
        \node[lbl] at ({2+\i*2},2.6) {$\alpha_{\a}$};
        \draw[arr] (h\i.north) -- ({2+\i*2},1.7);
    }

    % ---- output distributions ----
    \node[pout] (py1) at (5.0,2.0) {$p(y_1)$};
    \node[pout] (py2) at (7.0,2.0) {$p(y_2)$};
    \draw[arr] (h1.north) -- (py1.south);
    \draw[arr] (h2.north) -- (py2.south);

    % ---- inputs (context + previous word) ----
    \node[ctx] (z1) at (3.4,-1.5) {$z_1$};
    \node[yin] (y0) at (4.6,-1.5) {$y_0$};
    \node[ctx] (z2) at (5.4,-1.5) {$z_2$};
    \node[yin] (y1) at (6.6,-1.5) {$y_1$};
    \node[lbl, below=1mm of y0] {\textless START\textgreater};
    \foreach \n in {z1,y0} { \draw[arr] (\n.north) -- (h1.south); }
    \foreach \n in {z2,y1} { \draw[arr] (\n.north) -- (h2.south); }

    % ---- feedback: predicted word becomes next input ----
    \draw[->, thick, dlpurple] (py1.south) to[out=-90,in=90] (y1.north);
\end{tikzpicture}

    \caption{Show, Attend and Tell: at each step the decoder produces an attention map $\alpha_t$ over the CNN feature map, forming a context vector $z_t$ that, together with the previous word, drives the prediction $p(y_t)$.}
\end{figure}

The architecture of \textit{Show, Attend and Tell} consists of three main components:
\begin{itemize}
    \item \textbf{Encoder}: Similar to \textit{Show and Tell}, a CNN is used for feature extraction. 
    However, to preserve spatial information, the encoder avoids the fully connected layers that would otherwise collapse the spatial dimensions. 
    Instead, features are extracted directly from the final convolutional layer, producing a set of spatial feature vectors $\{a_1, a_2, \dots, a_L\}$. 
    Here, $L$ represents the number of regions in the image (e.g., $14 \times 14 = 196$), and each $a_i$ corresponds to the features of a specific region.
    
    \item \textbf{Attention Mechanism}: At each time step $t$, an attention network $f_{att}$ determines where the model should focus. 
    To do this, it computes a set of attention weights $\alpha_{t,i}$ for each region $i$, conditioned on the decoder's previous hidden state $h_{t-1}$ and the encoder's feature vectors $a_i$. 
    
    A \textbf{context vector} $z_t$ is then computed as a function of the attention weights $\alpha_{t,i}$ and the corresponding feature vectors $a_i$.
    The function could either be a weighted sum or a more complex non-linear transformation.
    This context vector dynamically captures the specific visual information relevant for generating the current word.
    
    \item \textbf{Decoder}: An LSTM generates the caption sequentially. 
    At time step $t$, the decoder receives the dynamic context vector $z_t$, the previous hidden state $h_{t-1}$, and the previously generated word $y_{t-1}$. 
    The forward pass of the LSTM updates the hidden state and produces a probability distribution over the vocabulary for the next word $y_t$. 
    By conditioning the generation on $z_t$, the decoder ensures that each word is grounded in the most relevant visual features of the image.
\end{itemize}


\subsection{Attention Types}

Attention mechanisms are generally categorized into two primary paradigms, distinguished by how they compute the context vector from the input features:

\begin{itemize}
    \item \textbf{Soft Attention}: This is a deterministic and fully differentiable mechanism. 
    It computes a continuous, weighted average of all input features (e.g., encoder hidden states) based on their respective attention weights. 
    Because the entire operation is continuous and differentiable, the model can be trained end-to-end using standard backpropagation.
    
    \item \textbf{Hard Attention}: This is a stochastic and non-differentiable mechanism. 
    Instead of a weighted average, it makes a discrete choice to focus on a single element (or a specific subset) of the input at each time step, effectively applying a binary mask over the features. 
    Since this discrete selection process breaks differentiability, the model cannot be trained via standard backpropagation; instead, it relies on reinforcement learning techniques (such as the REINFORCE algorithm) to estimate gradients by treating the attention selection as a stochastic policy.
\end{itemize}

To illustrate this distinction within the context of image captioning: \textit{soft attention} evaluates the entire image at every step, creating a blended representation by weighting all spatial regions, whereas \textit{hard attention} makes a discrete decision to "crop" and look exclusively at one specific region of the image, ignoring the rest entirely.

\subsection{Machine Translation}

We will now try to explore the attention mechanism in the context of machine translation, which is a more complex task than image captioning.
Again, we will first introduce the baseline model without attention, and then we will analyze how attention can be incorporated to improve performance.

\subsubsection{No Attention - Seq2Seq}
The \textbf{sequence-to-sequence} (Seq2Seq) model is a foundational architecture for machine translation. 
It consists of two main components:

\begin{itemize}
    \item \textbf{Encoder}: An RNN (such as LSTM or GRU) processes the input sequence, mapping it to a fixed-length context vector;
    
    \item \textbf{Decoder}: Another RNN generates the output sequence based on the context vector produced by the encoder.
\end{itemize}

This architecture first showed that deep neural networks could be effectively applied to end-to-end translation tasks,
achieving competitive performance with traditional statistical machine translation methods.

\paragraph{Problems}
Again, the Seq2Seq model suffers from the same information bottleneck as the image captioning model. In fact, it assumes that the 
last hidden state of the encoder can capture all the semantic information of the input sequence. 
This is particularly challenging, in particular for long sentences, since the representation of the input sequence is limited to a single vector.

\paragraph{Minor Improvements: RNN Encoder-Decoder}
An early refinement to the standard Seq2Seq architecture introduced a distinct \textbf{context vector}. Rather than directly passing the encoder's final hidden state to the decoder, the context vector is explicitly computed as a function of that hidden state.

Although this added flexibility yielded minor performance gains, it failed to resolve the fundamental information bottleneck. Because the context vector remains a fixed-length representation, the model inherently struggles to compress and retain all relevant information from longer input sequences.

\subsubsection{RNN with Attention}

Unlike earlier models that rely on a single, fixed context vector to summarize the entire input, the attention mechanism introduces a \textbf{dynamic context vector} computed at each decoder time step. 
This allows the network to adaptively "attend" to the most relevant parts of the input sequence as it generates each target word.

This architecture typically employs a Bi-directional RNN as the encoder. 
The forward and backward hidden states are concatenated to form a single, context-rich representation for each time step of the input sequence: $h_i = [\overrightarrow{h_i}; \overleftarrow{h_i}]$.

During the decoding step $t$, the model calculates an attention weight $\alpha_{t,i}$ for each encoder hidden state $h_i$. 
These weights are derived by scoring the alignment between the decoder's hidden state $s_t$ and each encoder state $h_i$.

The dynamic context vector $c_t$ is then constructed as the weighted sum of all encoder hidden states, scaled by their respective attention weights:
\[
c_t = \sum_{i=1}^{T_x} \alpha_{t,i} h_i
\]

\begin{figure}[H]
    \centering
    \begin{tikzpicture}[
    >=Stealth,
    enc/.style={rectangle, draw, fill=cyan!30, thick, minimum width=0.8cm, minimum height=0.45cm, font=\small},
    dec/.style={rectangle, draw, fill=red!45, thick, minimum width=0.7cm, minimum height=0.6cm, font=\small},
    sum/.style={circle, draw, thick, minimum size=0.55cm, inner sep=0pt},
    lbl/.style={font=\small},
    wlbl/.style={font=\footnotesize},
    arr/.style={->, thick}
]
    % ============ Bidirectional encoder ============
    \def\fy{0.65} \def\by{0.05}
    \foreach \i/\x in {1/1, 2/2.3, 3/3.6, T/5.4} {
        \node[enc] (f\i) at (\x,\fy) {$\overrightarrow{h}_{\i}$};
        \node[enc] (b\i) at (\x,\by) {$\overleftarrow{h}_{\i}$};
        \node[draw, dashed, gray, fit=(f\i)(b\i), inner sep=1.2mm] (col\i) {};
        \node[lbl] (x\i) at (\x,-0.95) {$x_{\i}$};
        \draw[arr] (x\i.north) -- (b\i.south);
    }
    % forward / backward recurrence
    \foreach \a/\b in {1/2, 2/3} {
        \draw[arr] (f\a.east) -- (f\b.west);
        \draw[arr] (b\b.west) -- (b\a.east);
    }
    \draw[arr, dashed] (f3.east) -- (fT.west);
    \draw[arr, dashed] (bT.west) -- (b3.east);
    \node at (4.5,0.35) {$\cdots$};

    % ============ Attention sum node ============
    \node[sum] (oplus) at (4.4,2.3) {$+$};

    % alignment weights from each column up to the sum
    \foreach \i/\w in {1/{t,1}, 2/{t,2}, 3/{t,3}, T/{t,T}} {
        \draw[arr] (f\i.north) -- node[wlbl, pos=0.22, fill=white, inner sep=0.5pt] {$\alpha_{\w}$} (oplus);
    }

    % ============ Decoder ============
    \node[dec] (sp) at (5.2,3.5) {$s_{t-1}$};
    \node[dec] (st) at (6.7,3.5) {$s_{t}$};
    \node[lbl] at (4.4,3.5) {$\cdots$};
    \node[lbl] at (7.5,3.5) {$\cdots$};
    \draw[arr] (4.7,3.5) -- (sp.west);
    \draw[arr] (sp.east) -- (st.west);
    \draw[arr] (st.east) -- (7.3,3.5);

    \node[lbl] (yp) at (5.2,4.5) {$y_{t-1}$};
    \node[lbl] (yt) at (6.7,4.5) {$y_{t}$};
    \draw[arr] (sp.north) -- (yp.south);
    \draw[arr] (st.north) -- (yt.south);

    % alignment model reads previous decoder state; context feeds new state
    \draw[arr] (sp.south) to[out=-90,in=60] (oplus);
    \draw[arr] (oplus) to[out=30,in=-120] node[lbl, pos=0.55, right=0.5mm] {$c_t$} (st.south);
\end{tikzpicture}

    \caption{RNN with attention: a bidirectional encoder produces hidden states $h_j$; alignment weights $\alpha_{t,j}$ combine them into the context vector $c_t$ that feeds the decoder state $s_t$ generating $y_t$.}
\end{figure}

\paragraph{Attention is Alignment}
When talking about machine translation, the attention mechanism can be interpreted as a form of \textbf{soft alignment} between the source and target sequences.
The attention weights $\alpha_{t,i}$ can be viewed as a probability distribution over the source tokens, indicating how much each source token contributes to the generation of the target token at time step $t$. 

\subsubsection{Local Attention}
Even though the attention mechanism we discussed alleviates the information bottleneck to some extent, it still suffers from the limitation of being a \textbf{global attention} mechanism.
In global attention, the decoder computes attention weights over all encoder hidden states at every time step, making it computationally expensive for long sequences.

\textbf{Local attention} addresses this by restricting the decoder to attend only to a small window of encoder states, rather than the whole sequence.
The idea is to design a mechanism that behaves similarly to \textit{hard attention} (focusing on a limited region) while remaining fully \textbf{differentiable}, and therefore trainable via standard backpropagation.

The core of the approach is, for each target time step $t$, to predict an \textbf{aligned position} $p_t$ in the source sequence: an estimate of which input region is most relevant for generating the current output word.
The context vector is then computed only from the encoder states inside a fixed-size window $[p_t - D, p_t + D]$ centered on $p_t$, where $D$ is a hyperparameter.

\begin{figure}[H]
    \centering
    \begin{tikzpicture}[
    >=Stealth,
    enc/.style={rectangle, draw, fill=attncontext!70, minimum width=0.5cm, minimum height=1.1cm, thick},
    dec/.style={rectangle, draw, fill=attnscore, minimum width=0.5cm, minimum height=1.1cm, thick},
    ctx/.style={rectangle, draw, fill=attncontext, minimum width=0.5cm, minimum height=1.2cm, thick},
    att/.style={rectangle, draw, fill=dlgray!40, minimum width=0.5cm, minimum height=1.2cm, thick},
    lbl/.style={font=\small},
    arr/.style={->, thick}
]
    % ---- Encoder row ----
    \foreach \i in {1,...,5} {
        \node[enc] (h\i) at (\i*1.15, 0) {};
        \draw[<-] (h\i.south) -- ++(0,-0.45);
    }
    \foreach \i [evaluate=\i as \j using int(\i+1)] in {1,...,4} {
        \draw[arr] (h\i.east) -- (h\j.west);
    }
    \node[lbl] at (3.45,-1.05) {$h_i$};

    % ---- Decoder cell ----
    \node[dec] (s) at (8.3,0) {};
    \draw[<-] (s.south) -- ++(0,-0.45);
    \draw[arr, dashed] (h5.east) -- (s.west);
    \node[lbl, right=1mm of s] {$s_t$};

    % ---- Attention window highlight ----
    \node[draw=dlred!80!black, very thick, rounded corners, fit=(h2)(h4), inner sep=2.5mm] (win) {};

    % ---- Context vector ----
    \node[ctx] (c) at (2.2, 2.6) {};
    \node[lbl, left=6mm of c] {$c_t$};
    \node[lbl] at (3.7,3.05) {Context vector};

    % ---- Local weights ----
    \node[att, minimum height=0.75cm, fill=dlgray!30] (a) at (4.0, 1.7) {};
    \node[lbl, above=0.5mm of a] {$a_t$};
    \node[lbl] at (5.4,1.3) {Local weights};

    % ---- Aligned position ----
    \node[lbl] (pt) at (6.5, 1.85) {$p_t$};
    \node[lbl] at (6.5,2.45) {Aligned position};

    % ---- Attentional hidden state ----
    \node[att] (sh) at (8.3, 2.6) {};
    \node[lbl, right=1mm of sh] {$\tilde{s}_t$};
    \node[lbl] at (5.8,3.9) {Attentional hidden state};
    \node[lbl] (yt) at (8.3, 4.05) {$y_t$};

    % ---- Attention layer box ----
    \node[draw, dash dot, thick, fit=(c)(a)(pt), inner sep=4mm,
          label={[lbl]above left:\textit{Attention Layer}}] (al) {};

    % ---- Arrows ----
    % encoder window -> context vector
    \foreach \i in {2,3,4} {
        \draw[arr] (h\i.north) -- (c.south);
    }
    % local weights feed the context vector
    \draw[arr] (a.west) -- (c.east);
    % context vector -> attentional hidden state
    \draw[arr] (c.east) -- (sh.west);
    % decoder state -> attentional hidden state
    \draw[arr] (s.north) -- (sh.south);
    % attentional hidden state -> output
    \draw[arr] (sh.north) -- (yt);
    % decoder state predicts the aligned position
    \draw[arr] (s.north) to[out=120,in=-20] (pt);
    % aligned position centers the window
    \draw[arr, dashed] (pt) to[out=210,in=20] (win.east);
\end{tikzpicture}

    \caption{Local attention with predictive alignment. The decoder state $s_t$ predicts the aligned position $p_t$, which centers a fixed window over the encoder states $h_i$. The context vector $c_t$ is formed from the window only and combined with $s_t$ into the attentional hidden state $\tilde{s}_t$ used to emit $y_t$.}
\end{figure}

In the \textbf{predictive alignment} variant, $p_t$ is computed from the decoder hidden state $s_t$ through a small learned network:
\[
p_t = T_x \cdot \mathrm{sigmoid}\!\left(v_p^\top \tanh(W_p s_t)\right)
\]
where $W_p$ and $v_p$ are learnable parameters.
Since $\mathrm{sigmoid}(\cdot) \in (0, 1)$, multiplying by the source length $T_x$ rescales the output into a valid position within the input sequence.

The window alignment is enforced by weighting the standard attention scores with a \textbf{Gaussian} centered on $p_t$:
\[
\alpha_{t,i} = \mathrm{softmax}(h_i^\top W_a s_t) \cdot \exp\!\left(-\frac{(i - p_t)^2}{2\sigma^2}\right)
\]
where $W_a$ is a learnable scoring matrix and the Gaussian standard deviation is typically set to $\sigma = D/2$.
The Gaussian term damps the contribution of encoder states far from $p_t$ toward zero, effectively concentrating attention around the predicted position.

\paragraph{Why is computing $p_t$ tricky?}
The natural idea, predicting a position and then directly selecting the encoder state at that index, is \textbf{not differentiable}, since indexing into a discrete position breaks the gradient flow (this is exactly the issue that hard attention faces).
The Gaussian weighting is what resolves this difficulty: because $p_t$ is a real-valued quantity that appears \textit{inside} the smooth exponential term, the alignment remains differentiable with respect to $p_t$, and therefore with respect to $W_p$ and $v_p$.
This allows the model to learn where to look through gradient descent, recovering the focused behavior of hard attention without sacrificing trainability.

\begin{figure}[H]
    \centering
    \begin{tikzpicture}[
    >=Stealth,
    enc/.style={rectangle, draw, fill=attncontext!70, minimum width=0.4cm, minimum height=0.9cm, thick},
    dec/.style={rectangle, draw, fill=attnscore, minimum width=0.4cm, minimum height=0.9cm, thick},
    ctx/.style={rectangle, draw, fill=attncontext, minimum width=0.4cm, minimum height=1.0cm, thick},
    att/.style={rectangle, draw, fill=dlgray!40, minimum width=0.4cm, minimum height=1.0cm, thick},
    lbl/.style={font=\footnotesize},
    arr/.style={->, thick}
]
% ===================== LOCAL =====================
\begin{scope}
    \foreach \i in {1,...,4} {
        \node[enc] (Le\i) at (\i*0.9, 0) {};
        \draw[<-] (Le\i.south) -- ++(0,-0.35);
    }
    \foreach \i [evaluate=\i as \j using int(\i+1)] in {1,...,3} {
        \draw[arr] (Le\i.east) -- (Le\j.west);
    }
    \node[dec] (Ls) at (4.5,0) {};
    \draw[arr, dashed] (Le4.east) -- (Ls.west);
    \node[lbl, below=0.5mm of Ls] {$s_t$};

    \node[draw=dlred!80!black, very thick, rounded corners, fit=(Le2)(Le3), inner sep=1.5mm] (Lwin) {};

    \node[ctx] (Lc) at (1.5,2.0) {};
    \node[lbl, left=4mm of Lc] {$c_t$};
    \node[att, minimum height=0.55cm, fill=dlgray!30] (La) at (2.7,1.25) {};
    \node[lbl, above=-0.5mm of La] {$a_t$};
    \node[lbl] (Lpt) at (3.6,1.4) {$p_t$};
    \node[att] (Lsh) at (4.5,2.0) {};
    \node[lbl] (Lyt) at (4.5,3.0) {$y_t$};

    \node[draw, dash dot, fit=(Lc)(La)(Lpt), inner sep=2.5mm] (Lal) {};

    \foreach \i in {2,3} { \draw[arr] (Le\i.north) -- (Lc.south); }
    \draw[arr] (La.west) -- (Lc.east);
    \draw[arr] (Lc.east) -- (Lsh.west);
    \draw[arr] (Ls.north) -- (Lsh.south);
    \draw[arr] (Lsh.north) -- (Lyt);
    \draw[arr] (Ls.north) to[out=130,in=-20] (Lpt);
    \draw[arr, dashed] (Lpt) to[out=215,in=25] (Lwin.east);

    \node[font=\small\bfseries] at (2.5,-1.15) {Local Attention};
\end{scope}
% ===================== GLOBAL =====================
\begin{scope}[xshift=7.6cm]
    \foreach \i in {1,...,4} {
        \node[enc] (Ge\i) at (\i*0.9, 0) {};
        \draw[<-] (Ge\i.south) -- ++(0,-0.35);
    }
    \foreach \i [evaluate=\i as \j using int(\i+1)] in {1,...,3} {
        \draw[arr] (Ge\i.east) -- (Ge\j.west);
    }
    \node[dec] (Gs) at (4.5,0) {};
    \draw[arr, dashed] (Ge4.east) -- (Gs.west);
    \node[lbl, below=0.5mm of Gs] {$s_t$};

    \node[ctx] (Gc) at (1.5,2.0) {};
    \node[lbl, left=4mm of Gc] {$c_t$};
    \node[att, minimum height=0.55cm, fill=dlgray!30] (Ga) at (3.0,1.25) {};
    \node[lbl, above=-0.5mm of Ga] {$a_t$};
    \node[att] (Gsh) at (4.5,2.0) {};
    \node[lbl] (Gyt) at (4.5,3.0) {$y_t$};

    \node[draw, dash dot, fit=(Gc)(Ga), inner sep=2.5mm] (Gal) {};

    \foreach \i in {1,2,3,4} { \draw[arr] (Ge\i.north) -- (Gc.south); }
    \draw[arr] (Ga.west) -- (Gc.east);
    \draw[arr] (Gc.east) -- (Gsh.west);
    \draw[arr] (Gs.north) -- (Gsh.south);
    \draw[arr] (Gsh.north) -- (Gyt);

    \node[font=\small\bfseries] at (2.5,-1.15) {Global Attention};
\end{scope}
\end{tikzpicture}

    \caption{Local vs global attention. Global attention scores all encoder states $h_i$ at every step, whereas local attention restricts the context to a window centered on the predicted position $p_t$.}
\end{figure}

